%% This file contains a number of tweaks that are typically applied to the main document.
%% They are not enabled by default, but can be enabled by uncommenting the relevant lines.

%%
%% Inline annotations; for predefined colors, refer to "dvipsnames" in the xcolor package:
%% https://tinyurl.com/overleaf-colors
%%
\pdfcompresslevel=9
\pdfobjcompresslevel=3

\newcommand{\red}[1]{{\color{red}#1}}
\newcommand{\blue}[1]{{\color{blue}#1}}
\newcommand{\orange}[1]{{\color{orange}#1}}
\newcommand{\magenta}[1]{{\color{magenta}#1}}
\newcommand{\todo}[1]{{\color{red}#1}}
\newcommand{\TODO}[1]{\textbf{\color{red}[TODO: #1]}}
\newcommand{\sz}[1]{\red{[SZ: #1]}}
\newcommand{\ig}[1]{\blue{[IG: #1]}}
\newcommand{\rf}[1]{\orange{[RF: #1]}}
\newcommand{\xs}[1]{\magenta{[XS: #1]}}
%%
%% disable for camera ready / submission by uncommenting these lines  
%%
% \renewcommand{\TODO}[1]{}
% \renewcommand{\todo}[1]{#1}

%%
%% work harder in optimizing text layout. Typically shrinks text by 1/6 of page, enable
%% it at the very end of the writing process, when you are just above the page limit
%%
\usepackage{overpic}
\usepackage{microtype}
\usepackage{graphicx} % Required for inserting images
\usepackage{tikz}
\usepackage[T1]{fontenc}
\usepackage{amsmath}
\usepackage{pifont}
\makeatletter
\let\oldtheequation\theequation
\renewcommand\tagform@[1]{\maketag@@@{\ignorespaces#1\unskip\@@italiccorr}}
\renewcommand\theequation{(\oldtheequation)}
\makeatother
\newcommand{\equationautorefname}{Eq.}

\usepackage[linesnumbered,ruled,vlined]{algorithm2e} % For algorithms
% \renewcommand{\algorithmcfname}{ALGORITHM}
\SetAlFnt{\small}
\SetAlCapFnt{\small}
\SetAlCapNameFnt{\small}
\SetAlCapHSkip{0pt}
\let\oldnl\nl% Store \nl in \oldnl
\newcommand{\nonl}{\renewcommand{\nl}{\let\nl\oldnl}}% Remove line number for one line
\definecolor{cmtClr}{HTML}{028A0F}
\newcommand{\mycmtfn}[1]{\textcolor{cmtClr}{\small #1}}

\usepackage{booktabs}
\usepackage{multirow}
\usepackage{threeparttable}
\usepackage{caption}
\usepackage{siunitx}

\usepackage{nicefrac}
\usepackage{comment}

\usepackage[outline]{contour}
\contourlength{0.075em}

\renewcommand{\algorithmautorefname}{Algorithm}
\newcommand{\sectionautorefname}{Section}
\newcommand{\subsectionautorefname}{Section}

\newlength{\resLen}

%%
%% fine-tune paragraph spacing
%%
\renewcommand{\paragraph}[1]{\vspace{.5em}\noindent\textbf{#1.}}

%%
%% globally adjusts space between figure and caption
%%
\setlength{\abovecaptionskip}{.5em}


%%
%% Allows "the use of \paper to refer to the project name"
%% with automatic management of space at the end of the word
%%
% \usepackage{xspace}
% \newcommand{\paper}{ProjectName\xspace}
\newcommand{\cmark}{\ding{51}}
\newcommand{\xmark}{\ding{55}}

%%
%% Commonly used math definitions
%%
% \DeclareMathOperator*{\argmin}{arg\,min}
% \DeclareMathOperator*{\argmax}{arg\,max}
\newcommand{\Real}{\mathbb{R}}
\newcommand{\sph}{\mathbb{S}^2}
\newcommand{\ex}{\mathbb{E}}
\newcommand{\pr}{\mathbb{P}}
\newcommand{\D}{\mathrm{d}}
\newcommand{\pluseq}{\mathrel{+}=}

\newcommand{\bmu}{\boldsymbol{\mu}}
\newcommand{\bSigma}{\boldsymbol{\Sigma}}
\newcommand{\Cest}{\langle C \rangle}
\newcommand{\loss}{\mathcal{L}}
\newcommand{\img}{\boldsymbol{I}}

\newcommand{\Mf}{M_\mathrm{f}}
\newcommand{\Mb}{M_\mathrm{b}}
\newcommand{\CestTot}{\Cest_\mathrm{tot}}

\newcommand{\brdf}{f_\mathrm{r}}
\newcommand{\vis}{T}
\newcommand{\Li}{L_\mathrm{in}}
\newcommand{\LiDir}{L_\mathrm{e}}
\newcommand{\LiInd}{\Li^\mathrm{ind}}

\newcommand{\bc}{\boldsymbol{c}}
\newcommand{\bx}{\boldsymbol{x}}
\newcommand{\bz}{\boldsymbol{z}}
\newcommand{\balpha}{\boldsymbol{\alpha}}
\newcommand{\bom}{\boldsymbol{\omega}}
\newcommand{\bomi}{\bom_\mathrm{in}}
\newcommand{\bomo}{\bom_\mathrm{out}}
\newcommand{\gShadow}{g^\mathrm{shadow}}
\newcommand{\alphaShadow}{\alpha^\mathrm{shadow}}

% \newcommand{\dCdbc}{\nicefrac{\partial C}{\partial\bc}}
\newcommand{\dCdbc}{\partial_{\bc} C}
\newcommand{\dCdbcEst}{\langle \dCdbc \rangle}
% \newcommand{\dCdbalpha}{\nicefrac{\partial C}{\partial\balpha}}
\newcommand{\dCdbalpha}{\partial_{\balpha} C}
\newcommand{\dCdbalphaEst}{\langle \dCdbalpha \rangle}

\newcommand{\imgwithtext}[3]{%
\begin{tikzpicture}
    \node[inner sep=0pt] (img) {\includegraphics[width=#1]{#3}};
    \node[anchor=south west, font=\scriptsize\bfseries, text=white] at (img.south west) {#2};
\end{tikzpicture}%
}

%%
%% Tigthen underline
%%
% \usepackage{soul}
% \setuldepth{foobar}